\documentclass[11pt,a4paper]{article}

% Tectonic = moteur XeTeX : UTF-8 natif, polices Latin Modern par defaut
\usepackage[french]{babel}
\usepackage[a4paper, top=2.2cm, bottom=2.0cm, left=2.2cm, right=2.2cm]{geometry}
\usepackage{xcolor}
\usepackage{titlesec}
\usepackage{enumitem}
\usepackage{booktabs}
\usepackage{tabularx}
\usepackage{tcolorbox}
\usepackage{fancyhdr}

\definecolor{accent}{HTML}{1F4E79}
\definecolor{accent2}{HTML}{2E7D32}
\definecolor{soft}{HTML}{EAF1F8}
\definecolor{linkblue}{HTML}{1565C0}

% Liens cliquables colores
\usepackage[colorlinks=true, linkcolor=accent, citecolor=accent2,
    urlcolor=linkblue, bookmarks=true, bookmarksnumbered=true, unicode=true,
    pdftitle={Mon ecosysteme numerique de these - Rugulopteryx okamurae},
    pdfauthor={Jihan El-khattabi},
    pdfsubject={Mini-projet Module 5 - Formation doctorale},
    pdfkeywords={ecosysteme numerique, ORCID, Zotero, PGD, Git, IA, Rugulopteryx okamurae}
]{hyperref}

% xurl : autorise le saut de ligne n'importe ou dans les URL (apres hyperref)
\usepackage{xurl}

% En-tete / pied de page
\pagestyle{fancy}
\fancyhf{}
\setlength{\headheight}{14pt}
\fancyhead[L]{\small\color{accent}\textbf{Mon écosystème numérique de thèse}}
\fancyhead[R]{\small\color{accent}Jihan El-khattabi -- LBECS}
\fancyfoot[C]{\small\thepage\,/\,\pageref{LastPage}}
\renewcommand{\headrulewidth}{0.4pt}
\renewcommand{\headrule}{\color{accent}\hrule width\headwidth height\headrulewidth}

% Style des sections
\titleformat{\section}{\color{accent}\large\bfseries}{\thesection.}{0.5em}{}
\titlespacing*{\section}{0pt}{10pt}{4pt}
\titleformat{\subsection}{\color{accent2}\normalsize\bfseries}{\thesubsection}{0.5em}{}
\setlist[itemize]{leftmargin=1.3em, itemsep=2pt, topsep=2pt}
\setlength{\parskip}{4pt}
\setlength{\parindent}{0pt}
\setlength{\emergencystretch}{2em}

% Boite de lien pour les livrables
\newtcolorbox{lienbox}[1]{%
  colback=soft, colframe=accent, boxrule=0.4pt, arc=2pt,
  left=6pt, right=6pt, top=3pt, bottom=3pt,
  before skip=4pt, after skip=6pt,
  fonttitle=\bfseries\color{accent}, title={#1}}

\begin{document}

% ===================== TITRE =====================
\begin{center}
{\color{accent}\Large\bfseries Mon écosystème numérique de thèse}\\[3pt]
{\large\itshape Mini-projet -- Module 5 «~Transformation numérique et recherche scientifique~»}\\[6pt]
{\color{accent2}\textbf{Thèse :} Étude phytochimique de \textit{Rugulopteryx okamurae} dans la Méditerranée marocaine}\\[4pt]
\rule{0.97\linewidth}{0.6pt}\\[3pt]
\small
\textbf{Doctorante :} Jihan El-khattabi \quad|\quad
\textbf{ORCID :} \href{https://orcid.org/0009-0002-7946-4934}{0009-0002-7946-4934}\\
\textbf{Laboratoire :} Biologie, Écologie et Santé (LBECS) \quad|\quad
\textbf{Date :} 15 novembre 2026
\end{center}
\vspace{2pt}

% ===================== SYNTHESE DES LIVRABLES =====================
\begin{lienbox}{Vue d'ensemble des livrables}
\renewcommand{\arraystretch}{1.15}
\begin{tabularx}{\linewidth}{@{}l l >{\raggedright\arraybackslash}X@{}}
\toprule
\textbf{N\textsuperscript{o}} & \textbf{Livrable} & \textbf{Lien} \\
\midrule
1 & Identité (ORCID) & \href{https://orcid.org/0009-0002-7946-4934}{orcid.org/0009-0002-7946-4934} \\
2 & Bibliothèque Zotero & \href{https://github.com/JihanElkhattabi/rugulopteryx-okamurae-these/blob/main/references.bib}{references.bib (22 réf.)} \\
3 & Plan de gestion des données & \href{https://github.com/JihanElkhattabi/rugulopteryx-okamurae-these/blob/main/docs/PGD_Rugulopteryx_okamurae.pdf}{PGD (2 pages, OPIDoR)} \\
4 & Dépôt Git & \href{https://github.com/JihanElkhattabi/rugulopteryx-okamurae-these}{github.com/JihanElkhattabi/rugulopteryx-okamurae-these} \\
5 & Charte IA & \href{https://github.com/JihanElkhattabi/rugulopteryx-okamurae-these/blob/main/Charte_IA.md}{Charte\_IA.md} \\
\bottomrule
\end{tabularx}
\end{lienbox}

% ===================== INTRODUCTION =====================
\section{Introduction}
La transformation numérique bouleverse les pratiques de la recherche scientifique et
impose à chaque doctorant de construire un \textbf{écosystème numérique} cohérent, où
chaque outil -- identifiant persistant, gestionnaire de références, plan de gestion
des données, contrôle de version et usage responsable de l'intelligence artificielle --
s'articule avec les autres pour garantir la traçabilité, la reproductibilité et
l'ouverture des résultats. Ce mini-projet, réalisé dans le cadre du Module~5 de la
formation doctorale, a précisément pour objectif de mettre en place cet écosystème
autour d'un sujet de thèse concret.

Ma thèse porte sur l'\textbf{étude phytochimique de \textit{Rugulopteryx okamurae}}
(E.Y.~Dawson) I.K.~Hwang, W.J.~Lee \& H.S.~Kim (Dictyotales, Ochrophyta), une algue
brune originaire du Pacifique Nord-Ouest dont l'invasion massive de la Méditerranée --
et particulièrement des côtes marocaines -- depuis 2015 menace la biodiversité
marine et les activités côtières. L'enjeu scientifique est double : caractériser le
profil métabolique de cette espèce (diterpénoïdes, phlorotannines, acides gras) et
évaluer le potentiel bioactif de ses extraits, dans une logique de valorisation de
cette biomasse invasive. Ce contexte impose une rigueur numérique particulière :
gestion fine de données analytiques volumineuses (spectres RMN, chromatogrammes),
bibliographie très mobilisée sur un sujet récent et en rapide expansion, et
transparence sur le recours aux outils d'IA générative.

Le présent rapport décrit les cinq livrables demandés (Slide~38), présente les
choix techniques retenus et fournit, pour chacun, un lien cliquable vers la
production correspondante.

% ===================== LIVRABLE 1 =====================
\section{Livrable 1 -- Identité : profil ORCID}
\begin{lienbox}{Lien du livrable}
\href{https://orcid.org/0009-0002-7946-4934}{https://orcid.org/0009-0002-7946-4934}
\end{lienbox}

L'identifiant ORCID (\textit{Open Researcher and Contributor ID}) constitue la
pierre angulaire de l'écosystème numérique : il associe de manière pérenne et
non ambiguë l'ensemble de mes productions scientifiques à mon identité, indépendamment
des homonymies et des changements d'affiliation. Mon profil
(\texttt{0009-0002-7946-4934}) a été créé puis complété dans l'interface
\url{https://orcid.org}. Il rassemble quatre rubriques essentielles :

\begin{itemize}
  \item \textbf{Biographie :} nom (Jihan El-khattabi), prénom, pays (Maroc) et mots-clés
        de recherche (phytochimie, macroalgues, espèces invasives).
  \item \textbf{Distinctions / éducation :} inscription en doctorat, formation
        universitaire et formations doctorales suivies (dont le Module~5).
  \item \textbf{Emploi :} affiliation actuelle au Laboratoire de Biologie, Écologie et
        Santé (LBECS).
  \item \textbf{Travaux :} section destinée à recevoir les publications, jeux de
        données et prépublications à mesure de leur parution, avec lien automatique
        vers CrossRef et DataCite.
\end{itemize}

La visibilité du profil ORCID est renforcée par son interopérabilité : il sera
systématiquement demandé lors des soumissions d'articles, des dépôts sur Zenodo/Nakala
et des demandes de financement. C'est l'identifiant unique qui relie entre eux tous
les autres composants de l'écosystème décrits ci-après.

% ===================== LIVRABLE 2 =====================
\section{Livrable 2 -- Bibliothèque Zotero}
\begin{lienbox}{Lien du livrable}
Export BibTeX (22 références) :
\href{https://github.com/JihanElkhattabi/rugulopteryx-okamurae-these/blob/main/references.bib}{references.bib}\\
Bibliothèque Zotero de groupe :
\href{https://www.zotero.org/groups/6764268/jihan_elkhattabi_bibliotheque_rugulopteryx}{\nolinkurl{zotero.org/groups/6764268/jihan_elkhattabi_bibliotheque_rugulopteryx}}
\end{lienbox}

J'ai constitué dans \textbf{Zotero} une bibliothèque de \textbf{22 références}
couvrant l'ensemble des dimensions du sujet, puis je l'ai exportée au format
\texttt{.bib} pour un usage direct avec LaTeX. Les références sont organisées par
\textbf{mots-clés thématiques}, ce qui permet un filtrage rapide dans l'interface
Zotero comme dans le fichier exporté :

\begin{itemize}
  \item \textit{Taxonomie \& systématique} -- Verlaque \textit{et al.} (2009), Hwang \textit{et al.} (2009) ;
  \item \textit{Écologie de l'invasion} (détroit de Gibraltar, Méditerranée) -- García-Gómez \textit{et al.} (2020), Altamirano \textit{et al.} (2021), Ruitton \textit{et al.} (2021), Laamraoui \textit{et al.} (2024) ;
  \item \textit{Invasion au Maroc} -- El~Aamri \textit{et al.} (2020), El~Madany \textit{et al.} (2024) ;
  \item \textit{Autres signalements méditerranéens} -- Bottalico \textit{et al.} (2024), Bellissimo \textit{et al.} (2024), Estévez \textit{et al.} (2022) ;
  \item \textit{Dynamique \& évaluation} -- Rosas-Guerrero \textit{et al.} (2025), Román \textit{et al.} (2025) ;
  \item \textit{Phytochimie \& bioactifs} -- Cuevas \textit{et al.} (2021, 2023), García-Cervantes \textit{et al.} (2025), Barcellos \textit{et al.} (2024) ;
  \item \textit{Valorisation \& bioraffinerie} -- De la Lama-Calvente \textit{et al.} (2024), Belhadj \textit{et al.} (2024), Trentin \textit{et al.} (2026) ;
  \item \textit{Gestion \& risques} -- Tsiriga \textit{et al.} (2024), Sempere-Valverde \textit{et al.} (2018).
\end{itemize}

\textbf{Vérification des rétractations.} Chaque référence a été contrôlée le
15/11/2026 via \href{https://www.crossref.org}{Crossref} (statut de l'identifiant DOI)
et la base \href{https://retractionwatch.com}{Retraction Watch}. Aucune rétractation
(\textit{retraction}) ni expression de préoccupation (\textit{expression of concern})
n'a été identifiée. Les métadonnées (auteurs, année, volume, DOI) ont été vérifiées
au moins au niveau du DOI pour les articles centraux (Cuevas \textit{et al.}, 2021 ;
Verlaque \textit{et al.}, 2009).

% ===================== LIVRABLE 3 =====================
\section{Livrable 3 -- Plan de gestion des données (PGD)}
\begin{lienbox}{Lien du livrable}
PGD (2 pages, modèle DMP OPIDoR) :
\href{https://github.com/JihanElkhattabi/rugulopteryx-okamurae-these/blob/main/docs/PGD_Rugulopteryx_okamurae.pdf}{docs/PGD\_Rugulopteryx\_okamurae.pdf}\\
Outil de rédaction : \href{https://dmp.opidor.fr}{dmp.opidor.fr} (DMPRoadmap)
\end{lienbox}

Le plan de gestion des données a été rédigé selon le \textbf{modèle DMP OPIDoR},
l'instance française de la plateforme DMPRoadmap recommandée par le HCERES. Il
s'articule autour de six sections qui couvrent l'intégralité du cycle de vie des
données de la thèse :

\begin{enumerate}[leftmargin=1.5em, itemsep=1pt]
  \item \textbf{Description des données} : trois familles -- échantillonnage
        (fiches terrain géoréférencées, photographies), analytiques (HPLC-DAD,
        HR-ESI-MS, RMN \textsuperscript{1}H/\textsuperscript{13}C, bioessais) et
        dérivées (tables, structures, figures) ; volume estimé 5--15~Go.
  \item \textbf{Normes \& métadonnées} : formats natifs conservés (FID, \texttt{.lcd},
        \texttt{.mzML}, \texttt{.csv}) ; standards Darwin Core (occurrences) et ISA-Tab
        (métabolomique) ; nomenclature normalisée des fichiers.
  \item \textbf{Stockage \& sécurité} : NAS du laboratoire avec sauvegarde nocturne,
        versionnement Git, chiffrement AES-256 des données sensibles, règle 3-2-1.
  \item \textbf{Partage \& dépôt} : Nakala (Huma-Num) et Zenodo avec DOI persistant,
        licence CC-BY-4.0, embargo maximal 12~mois pour les données brevetables.
  \item \textbf{Éthique \& légal} : respect du Code de la pêche maritime (loi 19-12),
        autorisations pour les aires protégées (parc national d'Al~Hoceima).
  \item \textbf{Responsabilités \& cycle de vie} : rôles doctorante / directeur /
        référent données ; révision annuelle ; archivage $\geq$~10~ans.
\end{enumerate}

Le PGD est joint au dépôt Git (dossier \texttt{docs/}) et sera importé dans OPIDoR
pour bénéficier de la version interactive et des modèles nationaux.

% ===================== LIVRABLE 4 =====================
\section{Livrable 4 -- Dépôt Git}
\begin{lienbox}{Lien du livrable}
Dépôt :
\href{https://github.com/JihanElkhattabi/rugulopteryx-okamurae-these}{github.com/JihanElkhattabi/rugulopteryx-okamurae-these}
\end{lienbox}

Le dépôt Git, désormais \textbf{en ligne} sur GitHub, héberge le socle versionné de
l'écosystème. Il a été structuré selon une \textbf{arborescence claire}, accompagnée
d'un fichier \texttt{README.md} décrivant l'objectif, la structure, les livrables et
la licence, d'un fichier \texttt{.gitignore} excluant les fichiers de compilation et
les données brutes volumineuses, et d'un fichier \texttt{CITATION.cff} fournissant les
métadonnées de citation (compatibles Zenodo) :

\begin{center}
\begin{minipage}{0.94\linewidth}
\small\ttfamily
rugulopteryx-okamurae-these/\\
\quad\textbar-- README.md \;\textperiodcentered\; LICENSE \;\textperiodcentered\; .gitignore \;\textperiodcentered\; CITATION.cff\\
\quad\textbar-- references.bib \;\textperiodcentered\; Charte\_IA.md\\
\quad\textbar-- docs/ \;\textbackslash\; PGD \;\textperiodcentered\; Rapport \;\textperiodcentered\; Guide\_GitHub.pdf\\
\quad\textbar-- data/ \;\textbackslash\; exemple\_metabolites.csv \;\textperiodcentered\; README.md\\
\quad\textbar-- scripts/ \;\textbackslash\; analyse\_phytochimique.py \;\textperiodcentered\; README.md\\
\quad\textbackslash-- results/ \;\textbackslash\; statistiques\_par\_classe.csv \;\textperiodcentered\; aire\_pic\_par\_compose.png \;\textperiodcentered\; README.md
\end{minipage}
\end{center}

La \textbf{licence} est double : CC-BY-4.0 pour les contenus (documents, données,
bibliographie) et MIT pour les scripts. Le dépôt contient un \textbf{script
d'analyse d'exemple} (\texttt{analyse\_phytochimique.py}), testé et fonctionnel, qui
lit un jeu de métabolites, calcule des statistiques par classe/site et génère une
figure. L'\textbf{historique} compte \textbf{huit commits} significatifs, traduisant
une démarche incrémentale et tracée :

\begin{itemize}
  \item \texttt{aeced97} -- Initialisation du dépôt (README, LICENSE, .gitignore) ;
  \item \texttt{2cfa4e4} -- Structuration de l'arborescence + \texttt{references.bib} (22 réf.) ;
  \item \texttt{7c79f29} -- Ajout de la charte personnelle d'usage de l'IA ;
  \item \texttt{17febc1} -- Ajout du PGD dans \texttt{docs/} ;
  \item \texttt{43bc509} -- Ajout du rapport final dans \texttt{docs/} + mise à jour du README ;
  \item \texttt{dcbb110} -- Fichiers d'exemple (data CSV, script, \texttt{results/}) + \texttt{CITATION.cff} ;
  \item \texttt{6cc3f22} -- README à jour (arborescence complète) + sorties d'exemple du script ;
  \item \texttt{907c2d9} -- Ajout du guide pas-à-pas de création du dépôt GitHub.
\end{itemize}

Le dépôt, désormais \textbf{poussé sur GitHub}, constitue la cible unique de
versionnement du code d'analyse phytochimique, des tables traitées et des scripts de
génération de figures, avec un \textbf{taggage de version} (\texttt{v0.1},
\texttt{v1.0}) à chaque jalon de la thèse et un archivage DOI via Zenodo.

% ===================== LIVRABLE 5 =====================
\section{Livrable 5 -- Test et charte IA}
\begin{lienbox}{Lien du livrable}
Charte personnelle d'usage de l'IA :
\href{https://github.com/JihanElkhattabi/rugulopteryx-okamurae-these/blob/main/Charte_IA.md}{Charte\_IA.md}
\end{lienbox}

\subsection*{Bilan de l'activité 2 (test d'un assistant IA)}
Lors de l'activité~2, j'ai soumis à un assistant conversationnel une requête sur les
composés bioactifs de \textit{R.~okamurae}. L'outil a produit un texte structuré et
plausible, mais l'examen critique a révélé trois limites majeures : (i)~\textbf{références
inventées} -- trois articles sur cinq n'existaient dans aucune base (Crossref, PubMed) ;
(ii)~\textbf{structures chimiques erronées} -- un diterpénoïde attribué à tort à la
classe des spatanes ; (iii)~\textbf{sources invérifiables} pour les dosages en
phlorotannines. Ce test confirme que l'IA est utile pour la formulation et la
synthèse, mais ne saurait remplacer la lecture critique des sources primaires.

\subsection*{Charte personnelle d'usage (résumé)}
Je m'engage à respecter six principes : (1)~\textbf{transparence} -- signaler tout
passage rédigé avec une IA, en nommant l'outil et la date ; (2)~\textbf{vérification
systématique} des références, chiffres et structures dans une source primaire ;
(3)~\textbf{aucune donnée fabriquée} -- retirer toute mention non contrôlée ;
(4)~\textbf{respect de la propriété intellectuelle} -- citer selon les normes en
vigueur ; (5)~\textbf{protection des données sensibles} -- ne jamais transmettre de
résultats non publiés à un outil tiers ; (6)~\textbf{responsabilité} -- rester seule
garante du contenu scientifique. Le texte intégral de la charte figure dans
\texttt{Charte\_IA.md}.

% ===================== CONCLUSION =====================
\section{Conclusion et perspectives}
L'écosystème numérique mis en place autour de ma thèse sur \textit{Rugulopteryx
okamurae} repose sur cinq briques interconnectées : un \textbf{identifiant ORCID}
qui me relie durablement à mes productions, une \textbf{bibliothèque Zotero} de
22~références vérifiées et exportables, un \textbf{plan de gestion des données}
conforme au modèle OPIDoR, un \textbf{dépôt Git} versionné et clairement structuré,
et une \textbf{charte d'usage de l'IA} traduisant un engagement responsable. Ces
outils ne fonctionnent pas en silos : l'ORCID alimente les dépôts, le PGD oriente
le versionnement Git, Zotero fournit la bibliographie du rapport, et la charte IA
encadre l'usage de tous les outils génératifs mobilisés. Cet écosystème évoluera
avec la thèse -- nouvelles références, nouveaux jeux de données, nouvelles versions
-- mais son squelette est désormais en place et reproductible.

En termes de \textbf{perspectives}, plusieurs actions prolongeront ce socle dès le
début du travail expérimental. La bibliothèque Zotero sera enrichie au rythme des
lectures, avec une veille automatisée sur les mots-clés «~\textit{Rugulopteryx}~»
et «~Dictyotales~» via les flux RSS de PubMed et Scopus. Le dépôt Git accueillera
les scripts d'analyse phytochimique (extraction chromatographique, traitement RMN,
figures), chaque jalon étant marqué par un tag de version. Le PGD sera révisé
annuellement et son contenu importé dans l'instance OPIDoR pour bénéficier des
modèles nationaux et de l'interopérabilité avec les entrepôts Nakala et Zenodo,
où les données analytiques seront déposées avec un DOI persistant. Enfin, la charte
IA sera actualisée à mesure que de nouveaux usages émergeront (assistance à la
rédaction, analyse de spectres), afin de maintenir un équilibre entre efficacité
et intégrité scientifique. Les références mobilisées (22~au total) sont listées,
organisées par mots-clés et vérifiées dans le fichier
\href{https://github.com/JihanElkhattabi/rugulopteryx-okamurae-these/blob/main/references.bib}{\texttt{references.bib}}
(livrable~2), auquel le présent rapport renvoie.

\label{LastPage}
\end{document}
